
\section{1. Introduction}

\subsection{1.1 Motivation}

Softmax attention computes an $L \times L$ score matrix, but the resulting 
distribution is \textbf{highly sparse}: the top few keys usually receive most 
of the attention mass. Empirical studies (Clark et al., 2019; Voita et al., 
2019; Zhang et al., 2023) consistently find:

\begin{itemize}
\item \textbf{Local dominance:} Most mass lies within a few tokens of the query.
\end{itemize}
\begin{itemize}
\item \textbf{A few global heads:} Some heads attend to distant content, typically
\end{itemize}
  across a small number of positions.
\begin{itemize}
\item \textbf{Content dependence:} Which distant positions matter varies with the
\end{itemize}
  input.

This suggests two orthogonal strategies:

1. \textbf{Fixed-window attention} (Longformer, Mistral): cheap, local.
2. \textbf{Content-based sparse attention} (Routing, Reformer): expensive, 
   content-aware.

\textbf{Neither alone is sufficient.} Fixed windows miss global dependencies; 
content-based methods require clustering or hashing overhead.

\subsection{1.2 Problem Statement}

Given a pretrained transformer with dense attention, we seek a modified 
module $\mathcal{A}'$ such that:

1. \textbf{Local capture:} Nearby tokens are always attended.
2. \textbf{Content-aware global:} A small number of distant tokens are selected 
   based on actual scores.
3. \textbf{Differentiable:} Selection is differentiable (or at least 
   trainable end-to-end).
4. \textbf{Vectorized:} No Python loop over queries.
5. \textbf{Quality:} $\Delta\text{Acc} \leq {{ACC_TOLERANCE}}\%$ on SST-2.
6. \textbf{Pretrained compatibility:} Only LoRA + $O(1)$ parameters are trainable.

\subsection{1.3 Our Approach: NSA}

\textbf{Native Sparse Attention} (NSA) combines a fixed sliding window and a 
dynamic top-$k$ selection. For each query $i$:

\[
\mathcal{S}_i = \mathcal{W}_i \cup \text{TopK}_{j \notin \mathcal{W}_i}(s_{ij})
\tag{1.1}
\]

where $\mathcal{W}_i = \{j : |i - j| \leq w\}$ is the local window and 
$\text{TopK}$ selects the $k$ highest-scoring keys outside the window.

\textbf{Key design choices:}

1. \textbf{No Python loop.} The top-$k$ selection uses \texttt{torch.topk} on a 
   masked score matrix. Fully vectorized.
2. \textbf{Union mask.} After top-$k$, we build a boolean mask and apply it 
   before softmax. No special kernels needed.
3. \textbf{Causal-aware top-$k$.} For decoders, the top-$k$ is computed 
   \textbf{after} causal masking, so it never picks future positions.
4. \textbf{Padding-safe.} Padding keys are excluded from top-$k$ selection.

\subsection{1.4 Contributions}

\begin{itemize}
\item \textbf{A hybrid sparse attention module.} NSA combines sliding windows with
\end{itemize}
  dynamic top-$k$ content-based routing.
\begin{itemize}
\item \textbf{Full vectorization.} No clustering, no hashing, no Python loop.
\end{itemize}
\begin{itemize}
\item \textbf{A production-hardened implementation.} Supports four architectures
\end{itemize}
  with RoPE, GQA, LoRA recovery, and safe injection.
\begin{itemize}
\item \textbf{A systematic evaluation.} Benchmarked on DistilBERT, RoBERTa, and GPT-2
\end{itemize}
  on SST-2 with ablations over window size, top-$k$, and causal handling.
\begin{itemize}
\item \textbf{Selection analysis.} We show that top-$k$ selected tokens vary with
\end{itemize}
  layer depth, revealing content-level head specialization.

\subsection{1.5 Paper Organization}

Section 2 reviews related work. Section 3 provides background on sparse 
attention. Section 4 presents NSA. Section 5 describes integration with four 
architectures. Section 6 details the experimental setup. Section 7 reports 
results. Section 8 analyzes the top-$k$ selection behavior. Section 9 
discusses practical guidance. Section 10 lists limitations. Section 11 
concludes.



\subsection{Notation}
